\begin{abstract}
Long-running molecular-simulation campaigns often require adaptive continuation, restart-dependent state progression, and provenance-aware control across multiple conditions. Although large language model (LLM) agents can support such workflows, continuous LLM involvement may add cost and nondeterminism to operations that are already governed by explicit rules. We developed Agent-MD, an architecture that combines rule-based campaign execution with event-triggered reasoning-agent or human review, in which supervisory reasoning constructs an executable campaign and a persistent rule-based campaign agent performs routine production. The framework was demonstrated using a path-dependent grand canonical Monte Carlo--molecular dynamics desorption campaign for montmorillonite. A high-level request was translated into five systems spanning exchangeable-cation and layer-charge effects and a relative-humidity path of $\mathrm{RH} = 0.9 \rightarrow 0.3 \rightarrow 0.1$, yielding 15 provenance-linked states. All states reached strict-pass archived status, while RH-local sampling varied from 2 to 42.1 million steps. The campaign comprised 414 campaign-agent operations and 120 run--analyze--decide cycles; only one production state required review, and no automated reasoning-agent invocation occurred during formal production. Analysis and frozen replay of an absolute-timestep interpretation error and a stale-artifact provenance conflict showed how structured evidence supported targeted correction and conservative recovery. The resulting hydration and basal-spacing trends were also physically consistent with known cation- and layer-charge effects. These results show that agentic scientific control need not rely on continuous LLM reasoning: routine actions can be selected and executed by a rule-based campaign agent, while flexible reasoning is reserved for decisions that could alter the scientific state or provenance of the campaign.
\end{abstract}
